%! AP Statistics — Unit 3, Lesson 3.4 — Homework KEY
\documentclass[10pt]{article}
\usepackage{apstats-article}
\usepackage{apstats-key}

\begin{document}

\pageheader{Unit 3, Lesson 3.4}{Homework — KEY}
\namedateperiod

\vspace{0.1in}

\begin{practicebox}
\small
\textbf{1.}
\begin{enumerate}[label=(\alph*), leftmargin=2em, itemsep=5pt]
  \item \ans{Voluntary response bias}
  \item \ans{Too high — only listeners with strong opinions (especially those who already favor later start times) call in. Neutral or opposing views are underrepresented.}
  \item \ans{Use a random sample of parents and students from school records rather than a call-in poll.}
\end{enumerate}

\textbf{2.}
\begin{enumerate}[label=(\alph*), leftmargin=2em, itemsep=5pt]
  \item \ans{Undercoverage bias}
  \item \ans{Households without landlines — primarily younger adults and lower-income households who rely on cell phones.}
  \item \ans{Younger adults tend to shop online more frequently; excluding them means the estimate of online shopping frequency is too low.}
\end{enumerate}

\textbf{3.}
\begin{enumerate}[label=(\alph*), leftmargin=2em, itemsep=5pt]
  \item \ans{Response bias (question-wording/leading question)}
  \item \ans{"Do you think the university should invest more in mental health resources?" or "How do you feel about university funding for mental health resources?"}
  \item Too \ans{high} because \ans{the phrase "don't you think" implies that agreeing is the expected response, pushing respondents to say yes.}
\end{enumerate}

\textbf{4.}
\begin{enumerate}[label=(\alph*), leftmargin=2em, itemsep=5pt]
  \item \ans{No — the initial selection was random.}
  \item \ans{Nonresponse bias}
  \item \ans{If those who complete the survey differ systematically from those who don't (e.g., more health-conscious, more engaged, more opinionated), the estimate will be distorted in that direction.}
\end{enumerate}
\end{practicebox}

\end{document}
